% GENERATED FILE. Edit canonical lessons, then run npm run generate:books.
% canonical-source-hash: fnv1a64:1c15e478ff7b97a1
% canonical-lessons: ES-C457-medicamento, ES-C457-encargar, ES-C457-guardia, ES-R457-repaso-guardia, ES-C457-sintesis-guardia

\chapter{La farmacia de guardia}
\label{ch:es-guardia}
\hlchaptermodality{457}

\begin{chapteropening}
\textbf{By the end of this chapter:} \emph{I can get a medicine that is not in stock: ask for it by the word on the box, ask them to order it in, find out when it arrives, and find the chemist that is open tonight.}

Read the message or notice this chapter's words exist for, then say the same thing back.
\end{chapteropening}

\section[el medicamento]{el medicamento — what the doctor ordered}

\label{lesson:ES-C457-medicamento}



\begin{quote}
Say \textbf{el médico} and \textbf{la farmacia}. One writes it down, the other hands it over.
\end{quote}

\subsection*{You'll want to know: el medicamento}
\textbf{el medicamento} — the medicine, the drug: the formal word, and the one on the packaging.

\emph{Un medicamento con receta} — a prescription medicine. \emph{Este medicamento no lo tenemos} — we don't have that one in.

In speech people say \textbf{la medicina} as readily, and for a single pill \textbf{la pastilla}. \emph{Medicamento} is what the box, the leaflet and the pharmacist will use.

\begin{grammarlens}[title={Grammar Lens: -mento and -miento are the same ending}]
You know \textbf{-miento} from \emph{el seguimiento}. \textbf{-mento} is the same Latin ending arriving by the other road:

\noindent
\begin{tabularx}{\linewidth}{@{}>{\raggedright\arraybackslash}X>{\raggedright\arraybackslash}X@{}}
\toprule
\textbf{worn down by speech} & \textbf{lifted from the page} \\
\midrule
el segui\textbf{miento} & el medica\textbf{mento} \\
el pensa\textbf{miento} & el docu\textbf{mento} \\
\bottomrule
\end{tabularx}

Both are masculine, so nothing is lost either way. What the spelling tells you is the \textbf{register}: a \emph{-miento} noun is usually everyday, a \emph{-mento} noun usually formal or technical. The pattern held for \emph{rueda} and \emph{rotativo}, and it holds here.
\end{grammarlens}

\begin{cousinweb}
Latin \textbf{medicus}, a doctor — and behind it \textbf{medēri}, to heal, to look after.

So \emph{el médico} and \emph{el medicamento} are one word: the healer, and the thing that does the healing. The family runs wide — \emph{la medicina}, \emph{remedio}, and English \textbf{remedy}, \textbf{medical}, \textbf{meditate}.

That last one is worth a second, and worth handling carefully. An older root does sit behind \emph{medeor}, and \textbf{modus}, \textbf{moderate} and \textbf{meditate} are independent evidence for it. What you should not do is gloss that root as \textquotedblleft{}measure\textquotedblright{} and then offer a doctor's judgement as the proof — that argues in a circle, as the lesson on \emph{el médico} has already said. The cousins are real; the story about why is not established.
\end{cousinweb}

\subsection*{Your turn}
Say these aloud:

\begin{itemize}
\raggedright
  \item a prescription medicine
  \item the pharmacy doesn't have that one
  \item what \emph{médico} and \emph{medicamento} share
\end{itemize}

\begin{tcolorbox}[breakable,colback=teal!4,colframe=teal!35!black,title={Before you move on}]
The medicine, on the box? (\textbf{El medicamento}.) Which ending is the formal one? (\textbf{-mento}.) What did \emph{medēri} mean? (\textbf{To heal, to look after}.)
\end{tcolorbox}
\section[encargar]{encargar — ask them to get it in}

\label{lesson:ES-C457-encargar}



\begin{quote}
Say \textbf{el cargo} and \textbf{el pedido}. Put \emph{en-} on the front of the first and you get a verb.
\end{quote}

\subsection*{You'll want to know: encargar}
\textbf{encargar} — to order something in, to ask for a thing to be got.

\emph{¿Me lo puede encargar?} — can you order it in for me? \emph{Lo tengo encargado} — it's on order. \emph{Encargar una tarta} — to order a cake.

With \textbf{-se de} it means to take charge of something: \emph{yo me encargo de la comida} — I'll see to the food. \emph{De eso se encarga la empresa} — the company handles that.

\begin{grammarlens}[title={Grammar Lens: encargar is not pedir}]
\noindent
\begin{tabularx}{\linewidth}{@{}>{\raggedright\arraybackslash}X>{\raggedright\arraybackslash}X@{}}
\toprule
\textbf{what is happening} & \textbf{the verb} \\
\midrule
asking for what is there & \textbf{pedir} \\
asking them to \textbf{get} what is not & \textbf{encargar} \\
\bottomrule
\end{tabularx}

In a restaurant you \emph{pides} a salad, because it is on the menu. At a chemist's you \emph{encargas} a medicine, because it has to be brought in. \emph{Pedir} takes something off a list; \emph{encargar} creates a job for somebody.

That is also why \emph{encargarse de} means taking responsibility: the load has been put on you.
\end{grammarlens}

\begin{cousinweb}
\textbf{en-} plus \textbf{cargar}, and \emph{cargar} is the \emph{carricare} from the lesson on \textbf{el cargo} — to load a cart.

So \emph{encargar} is to \textbf{load a task onto somebody}, and every sense follows from that one picture. Order a cake: the baker now carries it. Take charge: you are carrying it. \emph{El encargado} — the person in charge — is the one with the load on.

English has the same metaphor in \emph{charge}: you are \emph{in charge}, you \emph{take charge}, somebody is \emph{charged with} a duty. Two languages reached for a cart.
\end{cousinweb}

\subsection*{Your turn}
Say these aloud:

\begin{itemize}
\raggedright
  \item can you order the medicine in for me?
  \item I'll see to the food
  \item which verb you use for something on a menu
\end{itemize}

\begin{tcolorbox}[breakable,colback=teal!4,colframe=teal!35!black,title={Before you move on}]
To order something in? (\textbf{Encargar}.) To take charge of it? (\textbf{Encargarse de}.) What is being loaded, and onto what? (\textbf{A task, onto a person} — the cart again.)
\end{tcolorbox}
\section[la guardia]{la guardia — the one that stays open}

\label{lesson:ES-C457-guardia}



\begin{quote}
Say \textbf{guardar} and \textbf{la farmacia}. Somebody is keeping watch tonight, and it is one chemist in the town.
\end{quote}

\subsection*{You'll want to know: la guardia}
\textbf{la guardia} — duty, watch: the turn of being the one available.

\begin{quote}
\textbf{Farmacia de guardia.} Servicio nocturno.
\end{quote}

\emph{La farmacia de guardia} — the duty chemist, the one open when the rest are shut, posted on every pharmacy door. \emph{Estar de guardia} — to be on call. \emph{El médico de guardia} — the doctor on duty.

\textbf{El guardia} with the masculine article is a different thing: a police officer. The article carries the whole distinction.

\begin{grammarlens}[title={Grammar Lens: one word, two genders, two meanings}]
\noindent
\begin{tabularx}{\linewidth}{@{}>{\raggedright\arraybackslash}X>{\raggedright\arraybackslash}X@{}}
\toprule
\textbf{article} & \textbf{what it means} \\
\midrule
\textbf{la} guardia & the duty, the watch \\
\textbf{el} guardia & the officer keeping it \\
\bottomrule
\end{tabularx}

Spanish does this with a small set of words, and the rule inside it is consistent: the \textbf{feminine} is the abstract thing, the \textbf{masculine} is the person who does it. \emph{La guía} is the guidebook, \emph{el guía} is the person showing you round.

So \emph{la farmacia de guardia} is a chemist \textbf{on watch}, and \emph{un guardia en la puerta} is a man standing at it.
\end{grammarlens}

\begin{cousinweb}
Germanic \textbf{wardōn}, to watch over — the same root as \textbf{guardar}, which you already have for keeping and putting away.

Spanish took a shelf of Germanic words from the Visigoths and from Frankish, and they are recognisable by the \textbf{gu-} where English has plain \textbf{w-}:

\noindent
\begin{tabularx}{\linewidth}{@{}>{\raggedright\arraybackslash}X>{\raggedright\arraybackslash}X@{}}
\toprule
\textbf{Spanish} & \textbf{English} \\
\midrule
\textbf{gu}ardar & \textbf{w}ard \\
\textbf{gu}erra & \textbf{w}ar \\
\textbf{gu}ante & (glove, from the same Frankish \emph{want}) \\
\bottomrule
\end{tabularx}

English got \emph{guard} back from French, so \textbf{guard} and \textbf{ward} are the same Germanic word arriving twice — once directly, once round by way of Paris.
\end{cousinweb}

\subsection*{Your turn}
Say these aloud:

\begin{itemize}
\raggedright
  \item the duty chemist
  \item I'm on call tonight
  \item which article makes it a police officer
\end{itemize}

\begin{tcolorbox}[breakable,colback=teal!4,colframe=teal!35!black,title={Before you move on}]
The duty chemist? (\textbf{La farmacia de guardia}.) What does \emph{wardōn} mean? (\textbf{To watch over}.) What does Spanish \textbf{gu-} often answer to in English? (\textbf{W-} — \emph{guerra} and \emph{war}.)
\end{tcolorbox}
\section[(el medicamento, encargar, la guardia)]{(el medicamento, encargar, la guardia) — from cold}

\label{lesson:ES-R457-repaso-guardia}



\begin{quote}
Close the lessons before this one. Say all three: \emph{the medicine}, \emph{to order in}, \emph{the duty rota}.
\end{quote}

\begin{grammarlens}[title={Grammar Lens: the article doing the work}]
\noindent
\begin{tabularx}{\linewidth}{@{}>{\raggedright\arraybackslash}X>{\raggedright\arraybackslash}X@{}}
\toprule
\textbf{feminine — the thing} & \textbf{masculine — the person} \\
\midrule
la guardia — the watch & el guardia — the officer \\
la guía — the guidebook & el guía — the guide \\
la policía — the force & el policía — the officer \\
\bottomrule
\end{tabularx}

One pattern, and it is worth having in advance: the \textbf{feminine} names the abstract thing or the body, the \textbf{masculine} names a person doing it. A Spanish speaker hears the article and knows which you meant before the noun arrives.

So \emph{la farmacia de guardia} has nothing to do with a man on a door.
\end{grammarlens}

\begin{grammarlens}[title={Grammar Lens: two of the three were relatives}]
\noindent
\begin{tabularx}{\linewidth}{@{}>{\raggedright\arraybackslash}X>{\raggedright\arraybackslash}X@{}}
\toprule
\textbf{new word} & \textbf{what it was built on} \\
\midrule
\textbf{el medicamento} & \emph{el médico} — healer and the thing that heals \\
\textbf{encargar} & \emph{el cargo}, and \emph{el carrito} under that \\
\textbf{la guardia} & \emph{guardar}, from Germanic \emph{wardōn} \\
\bottomrule
\end{tabularx}

All three, in fact. \emph{Encargar} is the clearest: a cart, then a load, then a charge on a bill, then a task loaded onto somebody. Four lessons apart and one picture throughout.
\end{grammarlens}

\subsection*{Your turn}
Say these aloud:

\begin{itemize}
\raggedright
  \item can you order the medicine in?
  \item the duty chemist is open tonight
  \item what \emph{el guardia} means, with that article
\end{itemize}

\begin{tcolorbox}[breakable,colback=teal!4,colframe=teal!35!black,title={Before you move on}]
All three again, in Spanish, without looking. Then one sentence getting a medicine that is not in stock.
\end{tcolorbox}
\section[Practice]{(en la farmacia) — read it, then do it yourself}

\label{lesson:ES-C457-sintesis-guardia}



\begin{quote}
Say \textbf{la guardia} and \textbf{el medicamento}. The door below tells you where to find one of them after hours.
\end{quote}

\subsection*{Your turn}
A notice taped inside the glass, and the pharmacist's reply. Read both, then answer.

\textbf{1 — en la puerta}

\begin{quote}
Horario: 9:00 a 14:00 y 17:00 a 20:30.
\end{quote}

\begin{quote}
Fuera de ese horario, \textbf{farmacia de guardia} en la calle Mayor.
\end{quote}

\textbf{2 — en el mostrador}

\begin{quote}
Ese \textbf{medicamento} se nos ha terminado. Se lo \textbf{encargo} y lo tiene mañana por la tarde, si le parece.
\end{quote}

Say these aloud:

\begin{itemize}
\raggedright
  \item where to go at nine in the evening
  \item what the pharmacist offers to do
  \item when the medicine will be there
\end{itemize}

\begin{culture}
\textbf{Se lo encargo} is the sentence to keep whole. \emph{Encargo} rather than \emph{pido}, because the thing is not there to be asked for — it has to be brought in. And \emph{se lo} is two pronouns stacked: \emph{lo} is the medicine, \emph{se} is you. Spanish puts the person first and the thing second, which is the reverse of English.

\textbf{Si le parece} — \textquotedblleft{}if that suits you\textquotedblright{} — is how an offer is softened at a counter. It does the same work as the \emph{cómoda} of the delivery card, and both leave the decision with the customer without asking a question.

Notice too that the door says \textbf{fuera de ese horario} — the same \emph{fuera de} frame as \emph{fuera de servicio}, but here with a determiner. \emph{Fuera de servicio} is a fixed phrase and takes nothing; \emph{fuera de ese horario} points back at the hours printed two lines above, and pointing needs \emph{ese}.
\end{culture}

\begin{tcolorbox}[breakable,colback=teal!4,colframe=teal!35!black,title={Before you move on}]
Now your turn, out loud and then in writing. Four lines at a chemist's counter: ask for a medicine, hear that it has run out, ask them to order it in, and ask when it will be there. Then say where the duty chemist is.

Then check yourself: did you use \emph{encargar} rather than \emph{pedir}, and did the pronouns come out as \emph{se lo}?
\end{tcolorbox}
